% ============================================================
%  CHAPITRE 1 — À COMPLÉTER
%
%  Remplacez les textes entre [crochets] par votre contenu.
%  Pour ajouter un chapitre, copiez plutôt :
%      chapitres/modele-chapitre.tex
% ============================================================

\chapitre{1}{Titre du chapitre 1}

\begin{objectifs}
\begin{itemize}
  \item Comprendre [premier objectif du chapitre].
  \item Savoir [deuxième objectif du chapitre].
  \item Être capable de [troisième objectif du chapitre].
\end{itemize}
\end{objectifs}

\soussection{1.1}{Titre de la section 1.1}

Rédigez ici l'introduction, la description ou l'explication de votre
première section. Un paragraphe normal s'écrit directement dans le
fichier, sans commande particulière.

\soussoussection{1.1.1}{Titre de la sous-section 1.1.1}

Rédigez ici le premier point détaillé. Utilisez \code{un_nom_de_variable}
pour les éléments techniques et \motcle{un terme clé} pour attirer
l'attention sur un concept important.

\soussoussection{1.1.2}{Titre de la sous-section 1.1.2}

Ajoutez ici un deuxième point. Vous pouvez insérer un encadré, un bloc
de code, un tableau ou une image. Tous les exemples sont disponibles
dans le fichier \code{chapitres/demo.tex}.

\soussection{1.2}{Titre de la section 1.2}

Rédigez ici la deuxième partie du chapitre.

\begin{resume}
Résumez en quelques lignes les points essentiels du chapitre 1.
\end{resume}
